\subsection{A Galois normalization of the remaining mixed state}

\paragraph{D23.H2 (Galois normalization).}
Put $k=\Q^{\mathrm{ab}}$ and $\mathcal H_5=\GT\cap\ker\lambda$.
For every $h\in\mathcal H_5$ there is
$\sigma\in G_k\subset\mathcal H_5$ such that
\[
 \rho(h\sigma)=0\pmod3,\qquad
 a_2(h\sigma)=0,\qquad a_3(h\sigma)=\Delta(h).
\]
Consequently the existence of a full-GT element with $\Delta=1$
is equivalent to the existence of one with
$(\rho\bmod3,a_3,a_2)=(0,1,0)$.
This is a symmetry reduction of the same gate, not an exclusion
of any mixed state.

\begin{proof}
The supplied DELTA\_ON\_GT A.88 and A.102 identify
$\rho\bmod3=\Psi_2\bmod3$ on $\mathcal H_5$ and make this
an additive character. On $G_k$, the arithmetic comparison
in A.88 identifies $\Psi_2$ with the Kummer character of $2$,
up to the common sign determined by the monodromy marking.
Only its kernel and surjectivity modulo three are used here.
The source convention is recorded separately in A.100;
no equality between differently marked integral characters is
needed. The two frame bits are characters on the multiplier-one
gentle group by LM7\_PUSH E.65. Its Galois comparison and
DELTA\_ON\_GT A.12--A.13 give
\[
 a_2(\sigma)=a_3(\sigma)=\kappa_{r_0}(\sigma),
 \qquad r_0=(3/4)^{1/4}.
\]

The extension $k(\sqrt[3]{2})/k$ is cyclic of degree three.
Indeed $\zeta_3\in k$, and $\sqrt[3]{2}\notin k$: its real
cubic field is not normal over $\Q$, whereas every finite
subextension of the abelian extension $k/\Q$ is normal.
The extension $k(r_0)/k$ has degree two, as proved in the
supplied A.12 and A.25. Equivalently,
$k(r_0)=k(\sqrt[4]{3})$ because $\sqrt2\in k$.
These two finite Galois extensions have coprime degrees,
so their intersection is $k$ and their compositum has Galois
group $C_3\times C_2$. The restriction map from $G_k$ to that
group is surjective. Hence
\[
 G_k\longrightarrow\mathbf F_3\times\mathbf F_2,
 \qquad \sigma\longmapsto(\rho(\sigma),a_2(\sigma))
\]
is surjective, regardless of the common sign convention on
the first factor.

Choose $\sigma$ with
$\rho(\sigma)=-\rho(h)$ modulo three and
$a_2(\sigma)=a_2(h)$. Additivity of the three characters gives
the displayed normalization, since
$a_3(h)+a_2(h)=\Delta(h)$ on $\mathcal H_5$.
Multiplication by $\sigma$ takes place within actual full GT.
Thus this argument does not replace the lifting question by
an assertion about ordinary free-group words.
\end{proof}
